\textbf{Phần 1. Trắc nghiệm lựa chọn}
\setcounterpageref{ex}{0}

\begin{ex} %[4D5N25]
Công thức tìm số lớn khi biết tổng $T$ và hiệu $H$ của hai số là
\choice
{$(\text{Tổng} - \text{Hiệu}) : 2$}
{$\text{Tổng} + \text{Hiệu}$}
{$\text{Tổng} - \text{Hiệu}$}
{\True $(	ext{Tổng} + \text{Hiệu}) : 2$}
\loigiai{
Số lớn $= (\text{Tổng} + \text{Hiệu}) : 2$.\\
\textbf{Đáp án của câu hỏi là D}
}
\end{ex}

\begin{ex} %[4D5N25]
Công thức tìm số bé khi biết tổng $T$ và hiệu $H$ của hai số là
\choice
{\True $(	ext{Tổng} - \text{Hiệu}) : 2$}
{$(\text{Tổng} + \text{Hiệu}) : 2$}
{$\text{Tổng} - \text{Hiệu}$}
{$\text{Số lớn} + \text{Hiệu}$}
\loigiai{
Số bé $= (\text{Tổng} - \text{Hiệu}) : 2$.\\
\textbf{Đáp án của câu hỏi là A}
}
\end{ex}

\begin{ex} %[4D5H25]
Hai số có tổng bằng $70$ và hiệu bằng $20$. Số lớn là
\choice
{$40$}
{\True $45$}
{$50$}
{$35$}
\loigiai{
Số lớn $= (70 + 20) : 2 = 45$.\\
\textbf{Đáp án của câu hỏi là B}
}
\end{ex}

\begin{ex} %[4D5H25]
Hai số có tổng bằng $90$ và hiệu bằng $30$. Số bé là
\choice
{$25$}
{$35$}
{\True $30$}
{$60$}
\loigiai{
Số bé $= (90 - 30) : 2 = 30$.\\
\textbf{Đáp án của câu hỏi là C}
}
\end{ex}

\begin{ex} %[4D5H25]
Cả hai lớp 4A và 4B có tất cả $68$ học sinh. Lớp 4A có nhiều hơn lớp 4B là $4$ học sinh. Số học sinh của lớp 4A là
\choice
{$32$ học sinh}
{\True $36$ học sinh}
{$34$ học sinh}
{$38$ học sinh}
\loigiai{
Lớp 4A: $(68 + 4) : 2 = 36$ học sinh.\\
\textbf{Đáp án của câu hỏi là B}
}
\end{ex}

\begin{ex} %[4D5H25]
Một mảnh đất hình chữ nhật có nửa chu vi là $48\text{ m}$, chiều dài hơn chiều rộng $12\text{ m}$. Chiều dài của mảnh đất là
\choice
{\True $30\text{ m}$}
{$28\text{ m}$}
{$32\text{ m}$}
{$36\text{ m}$}
\loigiai{
Chiều dài $= (48 + 12) : 2 = 30\text{ m}$.\\
\textbf{Đáp án của câu hỏi là A}
}
\end{ex}

\begin{ex} %[4D5V25]
Một mảnh đất hình chữ nhật có chu vi $100\text{ m}$, chiều dài hơn chiều rộng $10\text{ m}$. Diện tích mảnh đất đó là
\choice
{$500\text{ m}^2$}
{$700\text{ m}^2$}
{\True $600\text{ m}^2$}
{$800\text{ m}^2$}
\loigiai{
Nửa chu vi: $100 : 2 = 50\text{ m}$.\\
Chiều dài: $(50 + 10) : 2 = 30\text{ m}$.\\
Chiều rộng: $50 - 30 = 20\text{ m}$.\\
Diện tích: $30 \times 20 = 600\text{ m}^2$.\\
\textbf{Đáp án của câu hỏi là C}
}
\end{ex}

\begin{ex} %[4D5V25]
Hai thùng chứa tất cả $150\text{ l}$ dầu. Nếu chuyển $15\text{ l}$ dầu từ thùng thứ nhất sang thùng thứ hai thì hai thùng có số lít dầu bằng nhau. Lúc đầu thùng thứ nhất chứa số lít dầu là
\choice
{$80\text{ l}$}
{$85\text{ l}$}
{$95\text{ l}$}
{\True $90\text{ l}$}
\loigiai{
Hiệu số dầu ban đầu: $15 \times 2 = 30\text{ l}$.\\
Thùng thứ nhất: $(150 + 30) : 2 = 90\text{ l}$.\\
\textbf{Đáp án của câu hỏi là D}
}
\end{ex}

\begin{ex} %[4D5V25]
Hai anh em có tất cả $45$ viên bi. Nếu anh cho em $5$ viên bi thì anh vẫn nhiều hơn em $3$ viên bi. Số bi của anh lúc đầu là
\choice
{\True $29$ viên}
{$27$ viên}
{$28$ viên}
{$30$ viên}
\loigiai{
Hiệu số bi lúc đầu: $5 \times 2 + 3 = 13$ viên.\\
Số bi của anh lúc đầu: $(45 + 13) : 2 = 29$ viên.\\
\textbf{Đáp án của câu hỏi là A}
}
\end{ex}

\begin{ex} %[4D5V25]
Tổng số tuổi của hai bố con hiện nay là $48$ tuổi. Bố hơn con $28$ tuổi. Tuổi của con hiện nay là
\choice
{$8$ tuổi}
{$12$ tuổi}
{\True $10$ tuổi}
{$14$ tuổi}
\loigiai{
Tuổi con: $(48 - 28) : 2 = 10$ tuổi.\\
\textbf{Đáp án của câu hỏi là C}
}
\end{ex}

\begin{ex} %[4D5V25]
Hai số chẵn liên tiếp có tổng bằng $206$. Số lớn là
\choice
{$102$}
{$106$}
{$108$}
{\True $104$}
\loigiai{
Hiệu hai số chẵn liên tiếp là 2.\\
Số lớn: $(206 + 2) : 2 = 104$.\\
\textbf{Đáp án của câu hỏi là D}
}
\end{ex}

\begin{ex} %[4D5C25]
Một trường tiểu học có $640$ học sinh. Số học sinh nam nhiều hơn số học sinh nữ là $40$ em. Sau đó trường nhận thêm $15$ học sinh nữ và $5$ học sinh nam. Lúc này số học sinh nam nhiều hơn số học sinh nữ là
\choice
{$20$ em}
{\True $30$ em}
{$35$ em}
{$40$ em}
\loigiai{
Hiệu mới: $40 + 5 - 15 = 30$ em.\\
\textbf{Đáp án của câu hỏi là B}
}
\end{ex}

\textbf{Phần 2. Trắc nghiệm đúng sai}
\setcounterpageref{ex}{0}

\begin{ex} %[4D5H25]
Hai số có tổng là $120$ và hiệu là $40$. Xét tính đúng sai:
\choiceTF
{\True Số lớn là $80$.}
{\True Số bé là $40$.}
{Số lớn gấp $3$ lần số bé.}
{\True Số lớn gấp $2$ lần số bé.}
\loigiai{
\begin{itemize}
    \item Số lớn: $(120 + 40) : 2 = 80$. Đúng.
    \item Số bé: $(120 - 40) : 2 = 40$. Đúng.
    \item $80 : 40 = 2$ (gấp 2 lần, không phải 3 lần). Ý 3 Sai, Ý 4 Đúng.
\end{itemize}
\begin{itemize}
    \item a) \textbf{Đúng}.
    \item b) \textbf{Đúng}.
    \item c) \textbf{Sai}.
    \item d) \textbf{Đúng}.
\end{itemize}
}
\end{ex}

\begin{ex} %[4D5V25]
Một thửa ruộng hình chữ nhật có chu vi $140\text{ m}$, chiều dài hơn chiều rộng $20\text{ m}$. Xét tính đúng sai:
\choiceTF
{\True Nửa chu vi thửa ruộng là $70\text{ m}$.}
{\True Chiều dài thửa ruộng là $45\text{ m}$.}
{Chiều rộng thửa ruộng là $35\text{ m}$.}
{\True Diện tích thửa ruộng là $1125\text{ m}^2$.}
\loigiai{
\begin{itemize}
    \item Nửa chu vi: $140 : 2 = 70\text{ m}$. Đúng.
    \item Chiều dài: $(70 + 20) : 2 = 45\text{ m}$. Đúng.
    \item Chiều rộng: $70 - 45 = 25\text{ m} \neq 35\text{ m}$. Sai.
    \item Diện tích: $45 \times 25 = 1125\text{ m}^2$. Đúng.
\end{itemize}
\begin{itemize}
    \item a) \textbf{Đúng}.
    \item b) \textbf{Đúng}.
    \item c) \textbf{Sai}.
    \item d) \textbf{Đúng}.
\end{itemize}
}
\end{ex}

\textbf{Phần 3. Trả lời ngắn}
\setcounterpageref{ex}{0}

\begin{ex} %[4D5H25]
Tổng hai số là $86$, hiệu hai số là $16$. Tìm số lớn.
\shortans{51}
\loigiai{
\textbf{Đáp số:} $51$\\
$(86 + 16) : 2 = 51$.
}
\end{ex}

\begin{ex} %[4D5H25]
Tổng hai số là $100$, hiệu hai số là $34$. Tìm số bé.
\shortans{33}
\loigiai{
\textbf{Đáp số:} $33$\\
$(100 - 34) : 2 = 33$.
}
\end{ex}

\begin{ex} %[4D5V25]
Hai bạn An và Bình có tất cả $54$ cái kẹo, An có nhiều hơn Bình $12$ cái kẹo. Tính số kẹo của An.
\shortans{33}
\loigiai{
\textbf{Đáp số:} $33$\\
$(54 + 12) : 2 = 33$.
}
\end{ex}

\begin{ex} %[4D5V25]
Hai số lẻ liên tiếp có tổng bằng $152$. Tìm số bé.
\shortans{75}
\loigiai{
\textbf{Đáp số:} $75$\\
$(152 - 2) : 2 = 75$.
}
\end{ex}

\begin{ex} %[4D5V25]
Tuổi bố và tuổi con cộng lại bằng $50$ tuổi, bố hơn con $26$ tuổi. Tính tuổi của bố.
\shortans{38}
\loigiai{
\textbf{Đáp số:} $38$\\
$(50 + 26) : 2 = 38$ tuổi.
}
\end{ex}

\begin{ex} %[4D5C25]
Một hình chữ nhật có chu vi $80\text{ cm}$. Nếu tăng chiều rộng thêm $5\text{ cm}$ và giảm chiều dài đi $5\text{ cm}$ thì hình chữ nhật trở thành hình vuông. Tính diện tích hình chữ nhật ban đầu theo xăng-ti-mét vuông.
\shortans{375}
\loigiai{
\textbf{Đáp số:} $375$\\
Nửa chu vi: $40\text{ cm}$.\\
Chiều dài hơn chiều rộng: $5 + 5 = 10\text{ cm}$.\\
Dài: $25\text{ cm}$, Rộng: $15\text{ cm}$.\\
Diện tích: $25 \times 15 = 375\text{ cm}^2$.
}
\end{ex}
